%============================================================
% SISTEMATIKA PENULISAN: PROGRES PA & BUKU PA (Bab 1-5)
% Di-input oleh chapters/bab1.tex via \input{\SistematikaFile}
% Dipanggil dari main_progres.tex dan main_buku.tex
%============================================================

\vspace{0.3cm}
\begin{description}[leftmargin=2.2cm, labelsep=0cm, itemsep=0.1cm]
  \item[{\makebox[2.2cm][l]{\textbf{Bab 1}}}] \textbf{Pendahuluan}\\
    Menjelaskan latar belakang, permasalahan, batasan masalah, tujuan, manfaat,
    dan sistematika penulisan.

  \item[{\makebox[2.2cm][l]{\textbf{Bab 2}}}] \textbf{Kajian Pustaka}\\
    Menjelaskan deskripsi permasalahan, teori penunjang, dan penelitian terkait
    yang mendukung penelitian proyek akhir.

  \item[{\makebox[2.2cm][l]{\textbf{Bab 3}}}] \textbf{Desain Sistem}\\
    Menjelaskan deskripsi solusi dan perancangan sistem yang diusulkan pada
    proyek akhir.

  \item[{\makebox[2.2cm][l]{\textbf{Bab 4}}}] \textbf{Eksperimen dan Analisis}\\
    Menjelaskan parameter eksperimen, karakteristik data, tempat dan waktu
    ujicoba, spesifikasi peralatan, hasil eksperimen, serta analisis hasil.

  \item[{\makebox[2.2cm][l]{\textbf{Bab 5}}}] \textbf{Penutup}\\
    Berisi kesimpulan terhadap penelitian yang dilakukan dan saran untuk
    pengembangan selanjutnya.
\end{description}
